%%%%%%%%%%%%%%%%%%%%%%%%%%%%%%%%%%%%%%%%%%%%%%%%%%%%%%%%%%%%%%%%%%
\section{High-order reconstruction of the solid on thin structures}\label{app:high_order_solid}
%%%%%%%%%%%%%%%%%%%%%%%%%%%%%%%%%%%%%%%%%%%%%%%%%%%%%%%%%%%%%%%%%%
% The old label of this material (Section 7.3, "Limitations") is kept as an alias
% so that cross-references from other sections still resolve; replace them by
% app:high_order_solid and then delete the alias.
\label{sec:lim_thin_solid}
% Moved from sectionLimitations.tex (Section 7.3) on 17 Nov 2026 (Philip's review).
% Data: paper-data/high_order_solid.csv. Runs: collapsible channel, local
% claude-runs/coll-ho697 (13 Nov 2026); FSI3, MeluXina fsi3ho697 (job 5447528).
% Condition numbers 5e17-1.4e19 on 54x3, 1.5e3-1.6e4 on 108x6 (solids4foam #698).

The velocity solids of this paper use the standard second-order cell-centred discretisation with linear reconstruction.
For thin, bending-dominated structures meshed with few, elongated cells across the thickness, this discretisation is too stiff: in a static test of the collapsible-channel wall it is between $16.7\%$ and $0.22\%$ too stiff over the solid meshes of the companion paper, whereas a high-order reconstruction is within $0.25\%$ on every mesh~\citep{CardiffFSIVerification2026}.
With the standard reconstruction, the accuracy of the coupled solution on such cases is therefore limited by the solid unless the solid is refined, which a non-conformal interface allows independently of the fluid (Section~\ref{sec:robustness_collapsible}).
This appendix compares the high-order residual of the geometrically nonlinear velocity solid (Section~\ref{sec:discrete_coupled_residual}), with a cubic reconstruction (order 3), and the standard reconstruction (order 1), within the quasi-monolithic coupling, on solid meshes with 2 to 32 cells across the thickness (Tables~\ref{tab:ho_collapsible} and~\ref{tab:ho_fsi3}).
These runs transfer the interface displacement to the fluid mesh through the points (Table~\ref{tab:mesh_motion_settings}), and their elapsed times come from runs sharing one machine, so they are indicative only.

\begin{table}[htbp]
	\centering
	\footnotesize
	\setlength{\tabcolsep}{3pt}
	\caption{Collapsible channel, fluid level 2, $\rho_S/\rho_F = 1$, $\Delta t = 1/160\,\mathrm{s}$, quasi-monolithic coupling: relative error of the first trough of the wall midpoint against oomph-lib, converged and at the same time step, Newton iterations per step and elapsed time (serial, runs sharing one Apple M1 Ultra workstation, indicative). Order 1: stabilisation of Equation~\eqref{eqn:rhie_chow_s} with $\alpha_d = 0.1$; order 3: jump stabilisation with factor $0.1$. On the $640\times8$ wall these runs give $-0.52\%$, against $-0.71\%$ in Section~\ref{sec:robustness_collapsible}; the difference comes from the interface displacement transfer and the assembly of the mesh-motion pseudo-solid (Section~\ref{sec:fluid_mesh_motion}).}
	\label{tab:ho_collapsible}
	\begin{tabular}{lrrrrrr}
		\toprule
		Order & Wall mesh & Cells & Error (\%) & Same $\Delta t$ (\%) & Newton & Elapsed (s) \\
		\midrule
		1 & $80\times2$    & 160    & $-12.34$ & $-12.23$ & 2.6 & 1379 \\
		1 & $160\times4$   & 640    & $-3.72$  & $-3.60$  & 2.8 & 1714 \\
		1 & $320\times8$   & 2560   & $-1.06$  & $-0.94$  & 2.7 & 1677 \\
		1 & $640\times8$   & 5120   & $-0.52$  & $-0.40$  & 2.7 & 1820 \\
		1 & $640\times16$  & 10240  & $-0.30$  & $-0.18$  & 2.8 & 1980 \\
		1 & $1280\times32$ & 40960  & $-0.11$  & $0.02$   & 2.7 & 2959 \\
		3 & $80\times2$    & 160    & \multicolumn{4}{l}{diverged at $t = 0.28\,\mathrm{s}$} \\
		3 & $160\times4$   & 640    & $-0.37$  & $-0.25$  & 2.8 & 2340 \\
		3 & $320\times8$   & 2560   & $-0.05$  & $0.08$   & 2.8 & 2572 \\
		\bottomrule
	\end{tabular}
\end{table}

\begin{table}[htbp]
	\centering
	\footnotesize
	\setlength{\tabcolsep}{3pt}
	\caption{Turek--Hron FSI3, fluid level 1 (Section~\ref{sec:fsi3_m}), quasi-monolithic coupling, flag meshes up to the level-1 flag ($430\times24$): relative errors (\%) against the Featflow level-4 values over $t = 5$--$6\,\mathrm{s}$. Both orders use the jump stabilisation with factor $0.1$, except order 3 on the $54\times3$ flag, which diverges with factor $0.1$ and is repeated with $1.0$. Eight runs of 16 processes side by side on one node took $32$--$36\,\mathrm{h}$ each for $t = 2$--$6\,\mathrm{s}$.}
	\label{tab:ho_fsi3}
	\begin{tabular}{lrrrrrr}
		\toprule
		Order & Flag mesh & Cells across & $u_x$ mean & $u_y$ ampl. & drag ampl. & lift ampl. \\
		\midrule
		1 & $54\times3$   & 3  & $-7.88$ & $-7.31$ & $25.54$ & $83.38$ \\
		1 & $108\times6$  & 6  & $0.80$  & $-1.39$ & $12.09$ & $32.29$ \\
		1 & $215\times12$ & 12 & $-0.43$ & $-0.20$ & $3.04$  & $10.05$ \\
		1 & $430\times24$ & 24 & $-0.40$ & $-0.07$ & $0.54$  & $4.04$ \\
		3 & $54\times3$   & 3  & $-9.65$ & $-5.59$ & $19.68$ & $56.44$ \\
		3 & $108\times6$  & 6  & $-0.08$ & $-0.22$ & $-0.31$ & $2.39$ \\
		3 & $215\times12$ & 12 & $-0.27$ & $-0.10$ & $-0.33$ & $1.91$ \\
		3 & $430\times24$ & 24 & $-0.29$ & $-0.08$ & $-0.36$ & $1.79$ \\
		\bottomrule
	\end{tabular}
\end{table}

On the collapsible channel the order-3 wall of $320\times8$ cells is within $0.08\%$ of oomph-lib at the same time step, closer than the order-1 wall of $640\times16$ cells and comparable with the order-1 wall of $1280\times32$, which has 16 times as many cells.
On the FSI3 flag, order 3 with six cells across the thickness is within $2.4\%$ in every quantity, and its lift-amplitude error changes only to $1.9$ and $1.8\%$ over the next two refinements, whereas order 1 approaches these values at about second order ($32.3$, $10.1$ and $4.0\%$); on the finest flag the two orders differ by $2.3\%$ of the reference value in the lift amplitude and by less than $1\%$ in the other quantities.
The Newton iterations per step are nearly unchanged by the order ($2.7$--$2.8$ on the collapsible channel).
The fluid dominates the cost of both cases, so the order and the solid mesh change the elapsed time little: order 3 costs $1.4$--$1.5$ times as much as order 1 on the same collapsible wall and about as much on the same FSI3 flag.
Its gain is that much coarser solid meshes reach a given accuracy: order 3 on the $108\times6$ flag is more accurate than order 1 on the $430\times24$ flag in the $u_x$ mean and the drag and lift amplitudes, in slightly less time, and order 3 on the $320\times8$ wall is about as accurate as order 1 on the $1280\times32$ wall in $0.87$ of its time.
With two or three cells across the thickness, the cubic fit through the cell centres cannot be determined across the thickness: the order-3 $80\times2$ wall stops at $t = 0.28\,\mathrm{s}$ when its Newton iterations rise from three to the iteration limit, and the $54\times3$ flag completes only with a ten-times larger stabilisation and still has errors of $6$--$56\%$.
In these tests order 3 therefore needs at least four cells across the thickness.
\todo{If Philip approves the demonstration runs: add the partitioned (Robin--Neumann or IQN-ILS) collapsible-channel results with the same order-3 walls, so that the high-order solid is shown with both couplings.}
